%!TEX root = ../report.tex
\documentclass[report.tex]{subfiles}
\begin{document}

\chapter{Solution}
\label{ch:solution}

\section{System Overview}
\label{sec:overview}

The system has three parts: a segmentation model that labels any frame
accurately but expensively, an optical-flow model that moves an existing mask
to the next frame cheaply, and a controller that chooses between them
(Figure~\ref{fig:system}).

\begin{figure}[t]
    \centering
    \begin{tikzpicture}
        % Spacing is measured between node BORDERS, never between centres, so
        % every connector keeps a fixed visible length whatever the height of
        % the boxes it joins. The chain is anchored on `flow` at the origin and
        % the frame pair is placed above it.
        \node[flowblk,  text width=46mm] (flow) at (0,0)
            {\textbf{SEA-RAFT}: forward $t{\to}t{+}1$\\
             \emph{and} backward $t{+}1{\to}t$};
        \node[flowblk,  text width=46mm, below=7mm of flow] (warp)
            {backward-warp the mask,\\then forward-backward check};
        \node[flowblk,  text width=46mm, below=7mm of warp] (cov)
            {coverage $\gamma_{t+1}$, AND-accumulated};
        \node[gate, text width=17mm, below=9mm of cov] (gate)
            {$\gamma_{t+1} \le \tau$?};
        \node[modelblk, text width=30mm, left=17mm of gate] (model)
            {\textbf{segmentation}\\\textbf{model}: re-infer,\\
             reset $\gamma \!\leftarrow\! 1$};
        \node[flowblk, text width=27mm, right=17mm of gate] (accept)
            {accept the\\warped mask};

        % The frame pair, 9mm of clear arrow above the dashed box.
        \coordinate (fin) at ($(flow.north)+(0,9mm)$);
        \node[plainblk, text width=20mm, anchor=south] (ft)
            at ($(fin)+(-23mm,7mm)$) {frame $I_t$};
        \node[plainblk, text width=20mm, anchor=south] (ft1)
            at ($(fin)+( 23mm,7mm)$) {frame $I_{t+1}$};

        % emit sits below whichever branch hangs lowest, on the centre line.
        \node[plainblk, text width=32mm, anchor=north] (emit)
            at ($(gate |- model.south)+(0,-15mm)$) {emit the mask for $I_{t+1}$};
        \coordinate (mrg) at ($(emit.north)+(0,7mm)$);

        \draw[seg] (ft.south)  |- (fin);
        \draw[seg] (ft1.south) |- (fin);
        \draw[fl]  (fin) -- (flow.north);
        \node[jn] at (fin) {};

        \draw[fl] (flow) -- (warp);
        \draw[fl] (warp) -- (cov);
        \draw[fl] (cov)  -- (gate);

        \draw[fl] (gate.west) -- node[lbl, above] {yes} (model.east);
        \draw[fl] (gate.east) -- node[lbl, above] {no}  (accept.west);

        \draw[seg] (model.south)  |- (mrg);
        \draw[seg] (accept.south) |- (mrg);
        \draw[fl]  (mrg) -- (emit.north);
        \node[jn] at (mrg) {};

        % Loop back for the next frame. Both ends must share one channel x, or
        % the "vertical" run is drawn diagonally: emit.west and ft.west do not
        % line up. The channel clears `model`, the widest block on this side.
        \coordinate (loopx) at ($(model.west)-(9mm,0)$);
        \draw[fl] (emit.west) -- (loopx |- emit.west)
                  -- (loopx |- ft.west) -- (ft.west);
        \node[lbl, rotate=90, anchor=south]
            at ($(loopx |- ft.west)!0.5!(loopx |- emit.west)-(1.5mm,0)$)
            {next frame};

        \begin{scope}[on background layer]
            \node[draw=cFlow, dashed, rounded corners=3pt, inner sep=3mm,
                  fit=(flow)(warp)(cov)] (flowbox) {};
        \end{scope}
        \node[lbl, anchor=north west, text width=30mm, align=left,
              text=cFlow!70!black] at ($(flowbox.north east)+(2mm,0)$)
            {paid on \emph{every} frame,\\before the decision\\can be made};
        % Above the model block, not beside it: to its left lies the loop-back
        % connector and its label.
        \node[lbl, anchor=south, text width=30mm, align=center, text=cModel]
            at ($(model.north)+(0,1.5mm)$) {paid on \PctModel\% of frames};
    \end{tikzpicture}
    \caption[Adaptive hybrid system overview]{The adaptive hybrid pipeline.
    Everything in the dashed box runs for each consecutive frame pair
    \emph{before} the threshold test, so its cost is unconditional; only the
    model call is conditional. That asymmetry, not the threshold, bounds the
    achievable speedup (Section~\ref{sec:sweep}).}
    \label{fig:system}
\end{figure}

\section{Segmentation Models}
\label{sec:segmodels}

Every model shares one shape: an encoder producing a four-level feature
pyramid at strides 4, 8, 16 and 32, and a UPerNet decoder~\cite{upernet2018}
fusing those levels into per-pixel logits, bilinearly upsampled to the input
resolution (Figure~\ref{fig:arch}). Every wrapper enforces the same contract,
input $(B, 3, H, W)$ and output $(B, C, H, W)$, which makes the backbones
interchangeable. Backbones start pretrained, from ImageNet-derived UPerNet
weights for ConvNeXt and Swin and from SAM~2~\cite{sam2_2024} for Hiera, and
only the six-class classifier trains from scratch, so convergence takes tens
rather than hundreds of epochs.

\begin{figure}[t]
    \centering
    \begin{tikzpicture}
        % Border-relative spacing: a fixed 7mm of visible arrow between boxes.
        \node[plainblk, text width=60mm] (in) at (0,0)
            {image $(B,3,H,W)$};
        \node[modelblk, text width=60mm, below=7mm of in] (enc)
            {\textbf{encoder}: ConvNeXt / Swin / Hiera / MiT};
        \node[plainblk, text width=60mm, below=7mm of enc] (pyr)
            {feature pyramid $C_1 \dots C_4$, strides 4, 8, 16, 32};
        \node[modelblk, text width=60mm, below=7mm of pyr] (ppm)
            {\textbf{PPM} on $C_4$: pool at
             $1{\times}1, 2{\times}2, 3{\times}3, 6{\times}6$,\\
             then convolve, upsample and concatenate};
        \node[modelblk, text width=60mm, below=7mm of ppm] (fpn)
            {\textbf{FPN} top-down: add into $C_3, C_2, C_1$,\\
             fuse all levels, $1{\times}1$ classifier};
        \node[plainblk, text width=60mm, below=7mm of fpn] (out)
            {logits $(B,C,H,W)$ after bilinear ${\times}4$};

        \draw[fl] (in)  -- (enc);
        \draw[fl] (enc) -- (pyr);
        \draw[fl] (pyr) -- (ppm);
        \draw[fl] (ppm) -- (fpn);
        \draw[fl] (fpn) -- (out);

        % lateral skip connections, routed down the left margin
        \draw[fl] (pyr.west) -- ($(pyr.west)-(10mm,0)$)
                  -- ($(fpn.west)-(10mm,0)$) -- (fpn.west);
        \node[lbl, rotate=90, anchor=south]
            at ($(pyr.west)!0.5!(fpn.west)-(11.5mm,0)$)
            {$C_1 \dots C_3$ laterals};
        % Balance the bounding box against the left-hand lateral routing.
        \path ($(in.east)+(14mm,0)$);
    \end{tikzpicture}
    \caption[Shared segmentation architecture]{The architecture shared by
    every trained model: an encoder, a four-level pyramid, and a UPerNet
    decoder combining pyramid pooling~\cite{pspnet2017} for global context with
    top-down FPN fusion for spatial detail.}
    \label{fig:arch}
\end{figure}


The Hiera variants~\cite{hiera2023} needed one fix to run at native
resolution. Hiera combines an interpolated global positional embedding with a
windowed embedding tiled by exact integer repetition. Training at $1024^2$
gives a $256 \times 256$ feature map, an exact multiple of the window size, so
the code worked; \SI{4}{K} input gives $960 \times 540$, and $540/8 = 67.5$ is
not integral, so it failed unconditionally. Replacing the exact tile with a
ceiling-tile-then-crop reproduces the old output wherever it was defined, so no
retraining was needed.

One training recipe serves every backbone, varying only batch size with memory
footprint (Table~\ref{tab:trainhp}). Random scale augmentation matters here
because UAVid is captured at roughly constant altitude and so contains little
natural scale variation.

\begin{table}[t]
    \centering
    \caption[Training hyper-parameters]{Training configuration, shared by every
    backbone except for batch size, which follows each model's memory
    footprint.}
    \label{tab:trainhp}
    \footnotesize
    \begin{tabular}{@{}ll@{}}
        \toprule
        \textbf{Setting} & \textbf{Value} \\
        \midrule
        Optimiser        & AdamW~\cite{adamw2019}, lr $10^{-4}$,
                           weight decay 0.05 \\
        Schedule         & cosine, with 2-epoch linear warmup \\
        Stopping         & early stop on a validation-mIoU plateau \\
        Input            & $1024^2$ random crops, \texttt{bfloat16} AMP \\
        Loss             & cross-entropy, ignoring class 255 \\
        Augmentation     & horizontal flip; random scale $0.75$ to $1.4$;
                           brightness \\
                         & and contrast; hue and saturation \\
        Batch size       & 8 (ConvNeXt-Large), 4 (Swin-Large and Hiera-Base+) \\
        \bottomrule
    \end{tabular}
\end{table}

\section{Optical Flow Backends}
\label{sec:flowbackends}

Both backends satisfy one interface: estimate flow for a pair, for a batch of
pairs, or for a batch returning a GPU-resident tensor, the last of which lets
the pipeline warp and score without a host round trip. Flow is returned in
pixels (Figure~\ref{fig:flowvec}).

\begin{figure}[t]
    \centering
    % Frame captions sit above the boxes and every correspondence arc is routed
    % below them. Arc labels carry an opaque background, so a label stays
    % legible even where an arc passes behind it.
    \begin{subfigure}[b]{0.48\textwidth}
        \centering
        \begin{tikzpicture}[scale=0.82,
                            arclbl/.style={lbl, fill=white, inner sep=1.6pt}]
            \draw[draw=cMuted, fill=cLight] (0,0) rectangle (2.4,1.4);
            \node[lbl, above] at (1.2,1.4) {source frame $t$};
            \draw[draw=cMuted, fill=cLight] (3.4,0) rectangle (5.8,1.4);
            \node[lbl, above] at (4.6,1.4) {target frame $t{+}1$};
            \fill[cModel] (1.50,0.75) circle (1.7pt);
            \node[lbl, above=0.3mm] at (1.50,0.77) {$q$};
            \fill[cFlow] (4.45,0.75) circle (1.7pt);
            \node[lbl, above=0.3mm] at (4.45,0.77) {$p$};
            \draw[-{Latex[length=1.7mm,width=1.4mm]}, cFlow, line width=0.8pt]
                (4.45,0.63) .. controls (4.1,-0.75) and (1.85,-0.75) .. (1.50,0.63);
            \node[arclbl, text=cFlow!70!black] at (2.98,-0.78)
                {$\mathbf{f}_{\text{bwd}}(p)$: where did $p$ come from?};
            % Match panel (b)'s depth so bottom alignment puts both frame
            % rows at the same height.
            \path (0,-1.45);
        \end{tikzpicture}
        \caption{backward warp: one source per target pixel}
    \end{subfigure}\hfill
    \begin{subfigure}[b]{0.48\textwidth}
        \centering
        \begin{tikzpicture}[scale=0.82,
                            arclbl/.style={lbl, fill=white, inner sep=1.6pt}]
            \draw[draw=cMuted, fill=cLight] (0,0) rectangle (2.4,1.4);
            \node[lbl, above] at (1.2,1.4) {source frame $t$};
            \draw[draw=cMuted, fill=cLight] (3.4,0) rectangle (5.8,1.4);
            \node[lbl, above] at (4.6,1.4) {target frame $t{+}1$};
            \fill[cModel] (1.50,0.75) circle (1.7pt);
            \node[lbl, above=0.3mm] at (1.50,0.77) {$q$};
            \fill[cFlow] (4.22,0.72) circle (1.7pt);
            \node[lbl, left=0.5mm] at (4.22,0.72) {$p$};
            \fill[cAccent] (4.98,0.98) circle (1.7pt);
            \node[lbl, right=0.5mm] at (4.98,0.98) {$r$};
            % step 1, a shallow arc below the frames
            \draw[-{Latex[length=1.7mm,width=1.4mm]}, cFlow, line width=0.8pt]
                (4.22,0.60) .. controls (3.9,-0.50) and (1.85,-0.50) .. (1.50,0.63);
            % step 2, a deeper arc so the two never touch
            \draw[-{Latex[length=1.7mm,width=1.4mm]}, cModel, line width=0.8pt]
                (1.50,0.57) .. controls (2.0,-1.50) and (4.9,-1.50) .. (4.98,0.86);
            % step 3, the residual, inside the target frame
            \draw[{Latex[length=1.3mm]}-{Latex[length=1.3mm]}, cAccent,
                  line width=1pt] (4.29,0.77) -- (4.91,0.93);
            \node[lbl, text=cAccent] at (4.52,1.18) {3.\ $\epsilon$};
            % Labels last, so their opaque backgrounds sit over the arcs.
            \node[arclbl, text=cFlow!70!black] at (2.87,-0.52)
                {1.\ $\mathbf{f}_{\text{bwd}}(p)$};
            \node[arclbl, text=cModel] at (3.40,-1.24)
                {2.\ $\mathbf{f}_{\text{fwd}}(q)$};
        \end{tikzpicture}
        \caption{FB check: valid iff $\epsilon < \SI{1.5}{px}$}
    \end{subfigure}
    \caption[Backward warp and forward-backward consistency]{The two
    propagation primitives. (a) Backward flow at target pixel $p$ names its
    source $q$, whose label is copied to $p$ by nearest-neighbour sampling.
    (b) A round trip through both flow directions returns to $r$; the distance
    $\epsilon$ from the original $p$ is the consistency error of
    Equation~\ref{eq:fb}.}
    \label{fig:warpfb}
\end{figure}

\begin{figure}[t]
    \centering
    \includegraphics[width=0.92\linewidth]{fig_deck_flowvectors.jpg}
    \caption[Estimated flow on a UAVid frame pair]{SEA-RAFT flow on the
    fastest-motion pair of a UAVid sequence, subsampled to one arrow per grid
    cell. Each arrow is the estimated displacement into the next frame, and
    labels are propagated along these vectors. The field is dominated by camera
    motion; the disagreement remaining at object boundaries is what the
    forward-backward check is meant to catch.}
    \label{fig:flowvec}
\end{figure}

SEA-RAFT~\cite{searaft2024} follows the RAFT recipe (Figure~\ref{fig:raft}):
encode features at $1/8$ resolution, build an all-pairs correlation volume,
regress an initial flow field, refine it for a fixed number of recurrent
iterations that look up the volume at the current estimate, then convex
upsample to full resolution. The configuration used is the \texttt{spring-L}
checkpoint with 12 refinement iterations, a ResNet-34 feature encoder,
correlation radius 4 and feature dimension 128.

\begin{figure}[t]
    \centering
    \begin{tikzpicture}
        % Border-relative spacing: a fixed 7mm of visible arrow between boxes.
        \node[plainblk, text width=56mm] (im) at (0,0) {$I_1, I_2$};
        \node[flowblk,  text width=56mm, below=7mm of im] (feat)
            {feature encoder, ResNet-FPN,\\at $1/8$ resolution};
        \node[flowblk,  text width=56mm, below=7mm of feat] (corr)
            {\textbf{all-pairs correlation volume}\\size $(H/8 \cdot W/8)^2$};
        \node[flowblk,  text width=56mm, below=7mm of corr] (gru)
            {regress an initial flow, then refine it:\\
             the GRU looks up the correlation};
        \node[gate, text width=15mm, below=8mm of gru] (done) {12 iters?};
        \node[flowblk,  text width=56mm, below=8mm of done] (out)
            {convex upsample ${\times}8$, giving the\\
             flow field $(H, W, 2)$ in pixels};

        \draw[fl] (im)   -- (feat);
        \draw[fl] (feat) -- (corr);
        \draw[fl] (corr) -- (gru);
        \draw[fl] (gru)  -- (done);
        \draw[fl] (done) -- node[lbl, right, pos=0.5] {yes} (out);

        % Refinement loop, routed up the right margin. Both ends share one
        % channel x, or the run is diagonal: the diamond is much narrower than
        % the block it returns to.
        \coordinate (loopx) at ($(gru.east)+(12mm,0)$);
        \draw[fl] (done.east) -- (loopx |- done.east)
                  -- (loopx |- gru.east) -- (gru.east);
        \node[lbl, anchor=west]
            at ($(loopx |- gru.east)!0.5!(loopx |- done.east)+(1.5mm,0)$) {no};

        \node[lbl, anchor=west, text width=30mm, align=left, text=cAccent]
            at ($(corr.east)+(3mm,0)$)
            {grows quadratically\\with resolution:\\the cost driver};
        % Balance the bounding box against the right-hand annotation.
        \path ($(im.west)-(34mm,0)$);
    \end{tikzpicture}
    \caption[SEA-RAFT iterative refinement]{SEA-RAFT's refinement
    loop~\cite{searaft2024,raft2020}. The correlation volume grows
    quadratically with resolution, so the internal working resolution is the
    dominant cost setting at \SI{4}{K}.}
    \label{fig:raft}
\end{figure}

The correlation volume dominates cost. For a frame of height $H$ and width
$W$ its size grows as $(H/8 \cdot W/8)^2$, so at \SI{4}{K} a single pair needs
roughly
\SI{20}{GB}, even though SEA-RAFT has far fewer \emph{parameters} than the
segmentation model. Parameter count is a poor proxy for cost here.

The remedy is a \texttt{scale} parameter, which runs the network at a lower
internal resolution and rescales the output vectors back to full-resolution
pixels. This project uses quarter resolution, \texttt{scale}${=}{-}2$. Flow
time falls by more than a factor of three, and a hard cuDNN limit that aborts
full-resolution batches whatever the free memory is avoided. The flow is still
returned at full input resolution, so mask resolution is unaffected, and
Section~\ref{sec:backends} measures the accuracy cost at under \num{0.01}
mIoU.

FlowNet2~\cite{flownet2_2017} is retained for comparison because SegProp uses
it. It is a single-pass cascade with no iterative refinement and no equivalent
resolution control, and it needs custom CUDA operators to run at all, which is
part of why SEA-RAFT is the default.

\section{Propagation Primitives}
\label{sec:primitives}

Three operations, applied to whole GPU-resident batches so that no per-frame
work returns to the host.

\subsection{Backward Warp}
\label{sec:warp}

To label target frame $t{+}1$ from a mask at $t$, the required flow is
\emph{backward}, from target to source. For each target pixel $(x, y)$,
$\mathbf{f}_{\text{bwd}}(x,y) = (u, v)$ says the content came from
$(x{+}u,\, y{+}v)$ in the source, so the target label is read from the source
mask at that location (Figure~\ref{fig:warpfb}a). Sampling is
\textbf{nearest-neighbour} rather than bilinear, since labels are categorical
and interpolating class 2 with class 5 yields either a meaningless value or a
valid but wrong class. Target pixels whose source falls outside the frame
become \texttt{ignore}. Backward warping is preferred to forward splatting
because it gives every target pixel exactly one source, leaving no holes and no
collisions.

\subsection{Forward-Backward Consistency}
\label{sec:fb}

Warped labels inherit the flow's errors. The forward-backward
check~\cite{sundaram2010,brox2004} detects unreliable correspondences by round
trip (Figure~\ref{fig:warpfb}b): step from target pixel $p$ to its claimed
source $q$ along backward flow, step from $q$ along forward flow sampled
bilinearly at that sub-pixel location, and measure the distance from $p$:
%
\begin{equation}
    \epsilon(p) = \bigl\lVert \mathbf{f}_{\text{bwd}}(p) + \mathbf{f}_{\text{fwd}}\bigl(p + \mathbf{f}_{\text{bwd}}(p)\bigr) \bigr\rVert_2 .
    \label{eq:fb}
\end{equation}
%
A pixel is valid when $\epsilon(p) < \theta$ and its source stayed in bounds,
with $\theta = \SI{1.5}{px}$ throughout.

\paragraph{What the check cannot see.} It tests whether the two flow fields
agree geometrically. It is blind to the nearest-neighbour rounding introduced
at every warp step, which accumulates at object boundaries while flow stays
self-consistent. This is why a frame can report high coverage and still carry
degraded labels on small moving objects (Section~\ref{sec:whereloss}).

\subsection{Coverage Accumulation}
\label{sec:coverage}

The check in Section~\ref{sec:fb} grades one warp. For the step from frame
$t$ to frame $t{+}1$ it marks every pixel true or false; call that map
$S_{t+1}$. On its own it says nothing about history: a pixel can pass this
step even though it was warped badly three frames ago, because the labels it
carries came from there.

The pipeline therefore keeps a second map, $V_t$, recording which pixels of
frame $t$ have passed \emph{every} step since the last model call. It is set
all-true immediately after inference. Updating it takes two operations. Since
$V_t$ is expressed in frame $t$'s pixel grid, it is first warped into frame
$t{+}1$ with the same backward flow used for the labels, then combined with
the new step's map by a per-pixel AND:
%
\begin{equation}
    V_{t+1} = \mathrm{warp}(V_t) \wedge S_{t+1},
    \qquad
    \gamma_{t+1} = \frac{1}{HW}\sum V_{t+1} .
    \label{eq:coverage}
\end{equation}
%
Coverage $\gamma_{t+1}$ is the fraction of the frame's $HW$ pixels that are
still true, and that single number is what the controller compares against
$\tau$. A pixel marked false can never become true again, so $\gamma$ only
falls while labels are propagated, and returns to 1 at the next model call.

\section{The Adaptive Algorithm and Implementation}
\label{sec:algorithm}

Algorithm~\ref{alg:adaptive} states the policy. Flow is computed for the whole
chunk at line~\ref{line:flowchunk}, before the threshold test at
line~\ref{line:gate}, so every frame pays for flow whether or not the model
then runs.

\begin{algorithm}[t]
    \caption{Adaptive hybrid inference}
    \label{alg:adaptive}
    \SetKwInOut{Input}{input}
    \SetKwInOut{Output}{output}
    \Input{frames $I_0..I_T$; model $f$; flow model $g$; threshold $\tau$; chunk size $K$}
    \Output{a dense mask $M_t$ for every frame}
    $M_0 \leftarrow f(I_0)$;\quad $V_0 \leftarrow \mathbf{1}$;\quad emit $(M_0, \text{``model''}, \gamma{=}1)$\;
    $s \leftarrow 0$\;
    \While{$s < T$}{
        $E \leftarrow \min(s+K,\, T)$\;
        $\{\mathbf{f}^{\text{fwd}}_i\}, \{\mathbf{f}^{\text{bwd}}_i\} \leftarrow g\bigl(\{(I_i, I_{i+1})\}_{i=s}^{E-1}\bigr)$ \tcp*{one batched launch}
        \label{line:flowchunk}
        \For{$i \leftarrow s$ \KwTo $E-1$}{
            $M' \leftarrow \mathrm{warp}(M_i,\, \mathbf{f}^{\text{bwd}}_i)$ \tcp*{nearest neighbour}
            $S \leftarrow \mathrm{fbcheck}(\mathbf{f}^{\text{fwd}}_i,\, \mathbf{f}^{\text{bwd}}_i,\, \theta)$\;
            $V' \leftarrow \mathrm{warp}(V_i,\, \mathbf{f}^{\text{bwd}}_i) \wedge S$;\quad
            $\gamma \leftarrow \mathrm{mean}(V')$\;
            \uIf{$\gamma \le \tau$}{ \label{line:gate}
                $M_{i+1} \leftarrow f(I_{i+1})$;\quad $V_{i+1} \leftarrow \mathbf{1}$\;
                emit $(M_{i+1}, \text{``model''}, \gamma{=}1)$\;
            }
            \Else{
                $M_{i+1} \leftarrow M'$;\quad $V_{i+1} \leftarrow V'$\;
                emit $(M_{i+1}, \text{``flow''}, \gamma)$\;
            }
        }
        discard chunk flow;\quad $s \leftarrow E$\;
    }
\end{algorithm}

Flow cannot be computed up front: the forward and backward fields for a
900-frame \SI{4}{K} video would occupy about \SI{120}{GB}. The pipeline
therefore streams, computing flow in windows of $K = 16$ steps and discarding
each window once its frames are labelled, so memory stays flat however long
the video is. Larger windows proved unreliable once a segmentation model
shared the GPU.

\end{document}
